%% notes-tex/database/database-expansion-hosts/sections/pk.tex
%% [[experiments.database.expansion-hosts]]
%% https://github.com/Mjvolk3/torchcell/tree/main/notes-tex/database/database-expansion-hosts/sections/pk.tex
%% SOURCE: hand-authored prose. Every count comes from the triage report
%% /scratch/projects/torchcell-scratch/tmp/deep-research/reports/Pk and next dataset wave triage.md
%% and from the two search passes recorded in
%% .../research_notes/Pk and next dataset wave triage/pichia_kudriavzevii.md, which carries the
%% per-finding verified, snippet and inferred labels this section preserves.

\section{Pichia kudriavzevii\secstatus{todo}}\label{sec:pk}

No published genome-wide essentiality or fitness screen exists for this species. There is no
Tn-seq, SATAY, Hermes, or piggyBac mutagenesis library, no genome-wide CRISPR knockout,
interference, or activation screen, and no arrayed deletion collection, for
\org{Pichia kudriavzevii}, \org{Issatchenkia orientalis}, \org{Candida krusei}, or
\org{Candida glycerinogenes} across 2010 to 2026. Douglass et al. 2018 established that those
four names are one species, with genomes collinear and 99.6 percent identical, so a search run
under all four names is the complete search for the host.

Each search axis is reported with its own counts, because a negative result is only citable
through the surface that produced it. A Europe PMC query crossing both current species names against ``CRISPR screen'',
``genome-wide screen'', ``essential genes'', ``transposon'', ``Tn-seq'', ``sgRNA library'',
``ribosome profiling'' and ``overexpression library'' returned 53 hits, of which 50 were
inspected at title level, with zero screens in this organism. A sweep of 74 preprints across
bioRxiv, medRxiv, Research Square, Preprints.org and Authorea found no screen, library, or
essentiality preprint. GEO DataSets queried by taxid 4909 returned 29 records and four series,
none of them a screen, and taxid 573 must not be used for the clinical name because it returns
1,328 unrelated records for a different genus. Lab-name searches established that Wheeldon's
genome-scale CRISPR platform exists for \org{Yarrowia lipolytica} and \org{Komagataella
phaffii} only, the latter as a six-fold-coverage library of 30,848 active guides covering more
than 99 percent of that genome, and has never been extended to this host. Essentiality in this
species exists only as output of the iIsor850 reconstruction over 850 genes, which the yeast
pool already excludes under the rule \texttt{not-a-dataset} because a reconstruction is not an
experiment.

The consequence for what the host can be used for is a limit on the label rather than on the
host. Gene-level fitness, essentiality, and genetic interaction are
unavailable as held-out labels in this species, so a transfer claim of the form trained on
yeast fitness and judged on \org{P. kudriavzevii} fitness has nothing to be judged against
today. What is available is isolate-level stress tolerance, antifungal susceptibility,
strain-level expression under acid stress, and per-locus integration fitness. Four axes remain
unsearched and are therefore unmeasured rather than negative: Chinese-language indices,
ArrayExpress and BioStudies, the Zenodo, Dryad, and figshare APIs, and the per-lab publication
pages of Alper, Blenner, and Jeffries.

\subsection{What does exist, led by the MAESTRO panel\secstatus{todo}}\label{sec:pk-have}

\notesec{MAESTRO, 170 isolates sequenced and 161 phenotyped} MAESTRO is the machine-learning
association method for small cohorts that the preprint is named for, and the panel it was run
on is whole-genome sequencing of 170
isolates with growth of 161 of them assessed across 57 stressors spanning heat, low pH,
organic acids, lignocellulosic inhibitors, and antifungals across azoles, polyenes, and
echinocandins, which is 9,177 growth measurements, inferred as the product of the two reported
axes. Copy-number variation in under 3.5 percent of genes has pleiotropic effects on stress
tolerance, a copy-number feature model reaches median $R^{2}$ 0.67 and median Pearson
$r$ 0.92 across the 57 conditions, and deletion of the tandem paralog pair
\gene{gene4260}/\gene{gene4261} produced effects across several stresses, which is the one
CRISPR-validated arm in the host. The panel is a preprint of 2025-09-04, bioRxiv
10.1101/2025.08.29.673034, corresponding author Ping-Hung Hsieh at LBNL. Its numbers conflict
across DOE pages, which report 162 strains, 305,435 SNPs, and 58 conditions against the
preprint's 170 strains and a project poster's 305,405 SNPs and 16,177 INDELs, so the figures
carried here are the preprint's and the disagreement is recorded rather than averaged. The
phenotype matrix itself has not been retrieved: bioRxiv returns 403 to every automated client
and Europe PMC's full-text route returns 500 for the preprint identifier.

\notesec{The low-pH expression atlas, 12 strains} Twelve isolates chosen for varying acid
tolerance, six tolerant and six susceptible, profiled at control pH near 5.5 and at pH 1.5,
sampled at 0 h and 1 h, with further arms testing 3-hydroxypropionate and citramalate, giving
roughly $2.4\times10^{5}$ expression measurements, inferred. It is the one omics resource in
the host built as a strain by condition matrix with a released processed count table, and the
table sits at Illinois Data Bank DOI 10.13012/B2IDB-1336177\_V1, which returns 403 to every
automated client.

\notesec{Douglass 2018, the cleanest public benchmark} Thirty-two sequenced isolates, 20
clinical and 12 environmental, with per-isolate MICs in mg/L for fluconazole, flucytosine,
amphotericin B, and micafungin, which is 128 MIC values read in full from the paper's own
table. Fluconazole MICs span 8 to 128 with a mode of 32. The same table gives ploidy, MAT
genotype, Illumina coverage, \gene{ERG11} and \gene{FKS1} variants, and copy status across the
\gene{ABC11}-\gene{ABC1} region per strain. Reads are at SRA SRP139299, the CBS573 annotated
reference at GenBank CP028773 to CP028778, and the second type strain at CP028531 to CP028535,
so this is the only row in the host that is fully public and fully provenanced today.

\notesec{The landing-pad screen, the only measured per-locus fitness} Forty-seven guide RNAs
across 21 intergenic regions, each characterized for cutting efficiency, integration
efficiency, cellular fitness, and GFP expression level, which is 188 readouts. Its artifacts
are a 23 KB package on a lab-hosted Box URL plus a GitHub repository, with no DOI and no
sequence accession, which is the retrieval risk Table~\ref{tab:asks} records.

\notesec{Four GEO series and the SD108 product line} GSE202655 holds six samples on L-malate
engineering and GSE159457 six samples on C4-dicarboxylate transporters in strain CY902 with raw
reads at SRP287255. GSE262836 and GSE282019 are two Delhi neonatal-intensive-care outbreak
series that together pair genome, expression, and fluconazole MIC across 10 RNA-seq samples.
The engineering line on SD108 covers succinate, 3-hydroxypropionate, citramalate,
glycolic acid, and xylose utilization, and its scale was overestimated on the first pass: Tran
et al. 2023 contributes about 10 distinct genotypes countable from the body text, with
Supplementary Table 5 unread and no sequence accession deposited at all, so the 100 to 250
genotype figure for the whole lineage is an upper-range estimate and is labeled one. The parts
panel of Cao et al. 2020 names 16 promoters and then ``and other tested promoters'', so its
total is unstated, and its deposit is one of at least eight Illinois Data Bank records for the
species, every one of which returned 403 on both passes.

\subsection{Four gene-identifier namespaces for one species\secstatus{todo}}\label{sec:pk-ids}

Four annotation namespaces are live for this one species, and the gap between the largest two
is wide enough that they are not inter-mappable by inspection. Table~\ref{tab:namespaces}
records each one with its gene count, the data that uses it, and the route that would join it
to the others.

The JGI MycoCosm assembly Issorie2, SD108 v2.0, is 21.69 Mbp over 130 scaffolds with 4,925
gene models, and 119 of those 130 scaffolds are predicted to be an alternate haplotype, with
separate primary and secondary allele gene-model tracks served. The whole engineering line
writes its loci in that namespace's \texttt{gNNNN} form. The MAESTRO preprint writes
\texttt{geneNNNN} instead, and whether \texttt{g4260} and \texttt{gene4260} are the same locus
is unverified, which is the kind of silent mismatch that corrupts a join without failing a
load. A third set is live in Ensembl Fungi for the same strain, assembly ASM76445v1, with 7,107
coding genes and 328 non-coding genes imported from INSDC. The clinical work uses a fourth,
the CBS573 reference annotated with YGAP by similarity and synteny to \org{S. cerevisiae},
about 10.8 Mb over five nuclear chromosomes.

The 4,925 against 7,107 difference is the reason a crosswalk has to be built from protein
sequence rather than from identifier similarity, and the haplotype split is the reason a naive
gene count in this species reads low against the 6,607 genes of the S288C reference. Cao et al.
2020 itself uses 10.8 Mbp as the genome size in its copy-number calculation, the haploid figure
rather than Issorie2's 21.69 Mbp, so two genome sizes already coexist in one lab's output. Under
the rule that the data is the data, a loader stores the released name verbatim in whichever
namespace the release used and records which namespace that is, and resolution across the four
is a downstream utility rather than a load-time gate.

%% GENERATED FILE -- do not hand-edit.
%% SOURCE: experiments/database/scripts/build_hosts_candidate_datasets_table.py
\begingroup
\footnotesize
\begin{longtable}{@{}L{22mm} L{34mm} L{20mm} L{30mm} L{66mm}@{}}
\caption[]{The gene-identifier namespaces this wave's released data speaks, per host.
\emph{Carries} is the model or gene count the namespace itself holds and \emph{Joins} is
what maps it to another namespace. One species in this wave has four live namespaces
whose model counts differ by more than two thousand genes, so a released gene id is not
a well-formed identifier until the namespace is named. A \emph{Joins} entry that says a
crosswalk has to be built is saying no join exists today: it is not a pointer to one.}
\label{tab:namespaces}\\
\toprule
Host & Namespace & Spelling & Carries & Joins \\
\midrule
\endfirsthead
\toprule
Host & Namespace & Spelling & Carries & Joins \\
\midrule
\endhead
\bottomrule
\endfoot
\org{P. kudriavzevii} & JGI Issorie2 gene models & \texttt{gNNNN} & 4,925 gene models & Nothing published. A crosswalk to the INSDC set has to be built by protein BLAST, and 119 of 130 scaffolds are predicted to be an alternate haplotype, so the gap is not only an identifier gap \\
\addlinespace[5pt]
\org{P. kudriavzevii} & MAESTRO gene spelling & \texttt{geneNNNN} & the same JGI model set under a different spelling & Assumed equal to the JGI spelling and not verified: g4260 must be checked against gene4260 rather than assumed equal, which is exactly the pair the MAESTRO CNV finding rests on \\
\addlinespace[5pt]
\org{P. kudriavzevii} & INSDC / Ensembl Fungi SD108 & \texttt{locus tag} & 7,107 coding genes & The 4,925-versus-7,107 gap is large enough that the two sets are not trivially inter-mappable. This is the set the genomes tier pins as pkudriavzevii\_SD108\_v1 \\
\addlinespace[5pt]
\org{P. kudriavzevii} & CBS573 / YGAP annotation & \texttt{YGAP id} & the clinical-isolate annotation & Used by the clinical work, including the MIC panel and the ERG11 and FKS1 variant calls, so a clinical row and an SD108 engineering row do not share a gene id at all \\
\addlinespace[5pt]
\org{P. putida} & KT2440 locus tags & \texttt{PP\_NNNN} & 5,729 genes in the served reference & Already served: the bacterial layer namespaces its five perturbation types by locus tag, so a PP\_ row writes the record type the schema holds \\
\addlinespace[5pt]
\org{P. putida} & Fitness Browser locusId & \texttt{orgId, locusId} & one row per gene per browser organism & refseq\_mapping.tsv maps org and locusId to a RefSeq accession and refseq\_assemblies maps orgId to an NCBI assembly, so the identifier bridge is supplied by the source rather than built \\
\addlinespace[5pt]
\org{P. aeruginosa} & PAO1 and PA14 reference numbers & \texttt{PA NNNN} & the two reference annotations the Tn-seq literature uses & GOLD\_84 carries PAO1 locus tags and GOLD\_115 is PA14 in LB, so the two gene lists are in different namespaces and must not be merged into one boolean \\
\addlinespace[5pt]
\org{P. aeruginosa} & Bakta de novo annotation & \texttt{per-isolate tag} & one annotation per newly sequenced isolate & Nothing maps a Bakta tag to a PA number, so an isolate survey annotated this way forces the alias layer rather than joining on a reference id \\
\addlinespace[5pt]
\org{S. cerevisiae} & SGD S288C R64-4-1 & \texttt{YALNNNC} & 6,607 genes in the shared vocabulary & The served namespace. A regulatory-DNA row is not gene-indexed at all: its identifier is the measured sequence, which is why those rows carry a reporter-locus or plasmid basis rather than a gene id \\
\addlinespace[5pt]
\end{longtable}
\endgroup


\subsection{The open leads\secstatus{todo}}\label{sec:pk-leads}

Three leads are unresolved rather than negative, and the first is time-sensitive. NCBI holds
713 BioProjects for this species and 357 of them match \texttt{resequencing}. Three sampled
ones, PRJNA1464946, PRJNA1464947 and PRJNA1464948, are titled as \org{Issatchenkia
orientalis} SD108 \texttt{IoSL10}, \texttt{IoSL16} and \texttt{IoSL19} resequencing,
registered on 2026-05-11 by the DOE Joint Genome Institute with no publication attached. A
series sequenced per isolate at that scale is either a mutagenized or an evolved SD108
library, and if it carries phenotypes it is the genome-wide resource this host currently
lacks. That is a hypothesis about an unopened deposit, not a finding.

The per-isolate deposition route for the MAESTRO panel is the second lead and is inferred
rather than verified. Three projects in the \texttt{IO\_NNN} naming exist, PRJNA571188 for
IO\_179, PRJNA571352 for IO\_048, and PRJNA583207 for IO\_044 gene expression profiling, each
registered as one BioProject per isolate rather than under an umbrella, which would explain
why no single accession for the panel appears in any search. Enumerating the complete
isolate-to-accession map needs either the preprint's data statement or a pass over all 357
summaries.

The third lead is a 69-condition expression matrix in this species, 5,140 genes by 69
conditions, described in a senior thesis from the Sinha group at UIUC that carries no
data-availability, code, deposition, or supplementary section of any kind. The provenance of
the expression and transcription-factor concentration data is one sentence that cites nothing.
So it is reachable by asking and by nothing else, and with a real provenance record a
69-condition expression matrix in this host would be worth more to the host's program than
anything else now public.
